\chapter{Codigo base sugerido para comecar}

Este capitulo junta um esqueleto minimo para iniciar sem perder sua estrutura.

\section{src/app.py}

\begin{lstlisting}[style=devbutterpython]
from pathlib import Path
import psycopg

from src.config.logger import configure_logging
from src.config.settings import Settings
from src.db.repository import HealthRecordRepository
from src.pipelines.sync_pipeline import run_sync_pipeline


def main() -> None:
    configure_logging()
    settings = Settings()

    xml_file = Path("data/export.xml").resolve()

    with psycopg.connect(settings.database_url) as connection:
        repository = HealthRecordRepository(connection)

        run_sync_pipeline(
            xml_file=xml_file,
            repository=repository,
            batch_size=settings.batch_size,
        )


if __name__ == "__main__":
    main()
\end{lstlisting}

\section{Comando para rodar}

\begin{lstlisting}[style=devbutterbash]
poetry run python src/app.py
\end{lstlisting}

\section{Dependencias iniciais}

\begin{lstlisting}[style=devbutterbash]
poetry add "psycopg[binary]"
poetry add --group dev pytest
\end{lstlisting}

Mais tarde:

\begin{lstlisting}[style=devbutterbash]
poetry add pandas streamlit plotly
\end{lstlisting}

\section{Queries iniciais de insight}

\subsection{Tipos mais comuns}

\begin{lstlisting}[style=devbuttersql]
SELECT
    type,
    COUNT(*) AS total
FROM health_records
GROUP BY type
ORDER BY total DESC;
\end{lstlisting}

\subsection{Batimento medio por dia}

\begin{lstlisting}[style=devbuttersql]
SELECT
    DATE(start_date) AS day,
    AVG(value_numeric) AS avg_heart_rate,
    MIN(value_numeric) AS min_heart_rate,
    MAX(value_numeric) AS max_heart_rate,
    COUNT(*) AS samples
FROM health_records
WHERE type = 'HKQuantityTypeIdentifierHeartRate'
GROUP BY DATE(start_date)
ORDER BY day;
\end{lstlisting}

\subsection{Passos por dia}

\begin{lstlisting}[style=devbuttersql]
SELECT
    DATE(start_date) AS day,
    SUM(value_numeric) AS steps
FROM health_records
WHERE type = 'HKQuantityTypeIdentifierStepCount'
GROUP BY DATE(start_date)
ORDER BY day;
\end{lstlisting}

\subsection{Media movel de 7 dias}

\begin{lstlisting}[style=devbuttersql]
WITH daily_steps AS (
    SELECT
        DATE(start_date) AS day,
        SUM(value_numeric) AS steps
    FROM health_records
    WHERE type = 'HKQuantityTypeIdentifierStepCount'
    GROUP BY DATE(start_date)
)
SELECT
    day,
    steps,
    AVG(steps) OVER (
        ORDER BY day
        ROWS BETWEEN 6 PRECEDING AND CURRENT ROW
    ) AS steps_7d_avg
FROM daily_steps
ORDER BY day;
\end{lstlisting}

\section{Comando de validacao com EXPLAIN}

\begin{lstlisting}[style=devbuttersql]
EXPLAIN ANALYZE
SELECT
    DATE(start_date) AS day,
    AVG(value_numeric) AS avg_heart_rate
FROM health_records
WHERE type = 'HKQuantityTypeIdentifierHeartRate'
GROUP BY DATE(start_date)
ORDER BY day;
\end{lstlisting}

Rode antes e depois dos indices para observar o plano de execucao.

\section{Proxima tarefa pratica}

A proxima tarefa deve ser implementar a Milestone 1:

\begin{enumerate}
  \item criar \code{docker-compose.yml};
  \item criar tabelas;
  \item criar \code{Settings};
  \item criar \code{configure_logging};
  \item criar \code{HealthRecord};
  \item criar testes de \code{clean_tag} e \code{safe_float};
  \item rodar um parse pequeno com fixture.
\end{enumerate}

So depois disso processe o arquivo gigante.
